%%%%%%%%%%%%%%%%%%%%%%%%%%%%%%%%%%%%%%%%%%%%%%%%%%%%%%%%%%%%%%%%
%General Packages (alphabetic order):						   %
%%%%%%%%%%%%%%%%%%%%%%%%%%%%%%%%%%%%%%%%%%%%%%%%%%%%%%%%%%%%%%%%
\usepackage[acronym]{glossaries} 	% For acronyms and gloassaries
\usepackage{adjustbox}              % For adjusting the box size (of tables)
\usepackage{amsmath}
%\usepackage{algorithmic}
\usepackage{cite}					% either cite or natbib
\usepackage{hyperref}                    %clickable references      % load after most stuff
\usepackage{cleveref}                    %\Cref \cref               % load after hyperref
\usepackage[width=160mm,
bindingoffset=14mm,
top=24mm,
bottom=27mm]{geometry}       % proper geometry for binding % load after hyperref

\hypersetup{
	colorlinks=true,
	linkcolor=blue,%linkcolor={blue!50!black},
	citecolor=blue,%red,%citecolor={blue!50!black},
	urlcolor=blue%urlcolor={blue!80!black}
}
%\usepackage{enumerate} 			% for alphanumerical enumerate  [(a)] [(i)]. Commented in favor of enumitem, that one seems more complete and has more support/features.
\let\labelindent\relax
\usepackage{enumitem}				% fot no indentation in itemize and roman or alphanumerical enumerate [(a)] [(i)]
\usepackage[T1]{fontenc}			% font
\usepackage{geometry}				% for setting page margin
\usepackage{graphicx}
	\graphicspath{{./figures/}}		% to set the figures path
%\usepackage[hidelinks]{hyperref}	% auto referencing figures, equations, etc., cleverref is a better package however
\usepackage[utf8]{inputenc}			% font
\usepackage{libertine}				% to use the libertine font
\usepackage{multirow}				% multirow for tables
\usepackage{tabularx} 				% extra features for tabular environment
\usepackage{threeparttable}
\usepackage{todonotes} 				% to use comments and to-do notes
%some useful options:				[disable, prependcaption, textsize=footnotesize]{todonotes} 
\usepackage{setspace}				% for setting line space to 1.5
\usepackage{subfigure}
\usepackage{soul} 					% for highlighting
\usepackage{url}					% for URLs and internet links
\usepackage{xargs}    				% for newcommandx definitions e.g. notes & colors
%\usepackage[svgnames]{xcolor}		% is already loaded by the document class









